one can replace the @@M0@@ by the @@M1@@ in this diagram, and obtain a diagram of the desired kind in terms of the Chow rings of @@M2@@ and @@M3@@.

\sgaStableTarget{sga6.e10.rem.7-13-13}\textbf{Remark 7.13.13.} Let @@M4@@ be a smooth projective scheme over an algebraically closed field. It follows trivially from the formalism developed in (SGA 5 IV) that, if @@M5@@ is a cycle of codimension @@M6@@ on @@M7@@ whose image in @@M8@@ is zero, then @@M9@@ is numerically equivalent to zero. One of the standard conjectures in the theory of algebraic cycles asserts that the converse is also true. Moreover, by (VII 4.11), it is immediate that one obtains an equivalent conjecture by replacing the group of algebraic cycles with the group @@M10@@, using definition (7.13.8), or again with the group @@M11@@, by stating that numerical equivalence to zero of an element of @@M12@@ is equivalent to vanishing of its augmentation and of its Chern classes in the @@M13@@.

With this understood, it follows immediately from the commutativity of (7.13.10) that, if @@M14@@ satisfies the preceding conjecture, then the specialization homomorphism @@M15@@ (and, a fortiori, @@M16@@) is compatible with numerical equivalence, as conjectured in 7.9. Note also that the conjecture invoked for this purpose (numerical equivalence implies @@M17@@-cohomological equivalence) is weaker than the one considered in 7.12.5.

\sgaStableTarget{sga6.e10.sec.7.14}\subsection*{7.14. Specialization homomorphism for groups of cycles and cycle classes.}
If @@M18@@ is a noetherian scheme, denote by @@M19@@ (resp. @@M20@@) the group of algebraic cycles on @@M21@@ purely of dimension @@M22@@ (resp. purely of codimension @@M23@@). When @@M24@@ is of finite type, one then defines
